% =============================================================================
% MLSys·im Tutorial Tutorial — Parts 0–4 (Morning Session)
% =============================================================================
\documentclass[aspectratio=169, 12pt]{beamer}
\usepackage{../../../slides/assets/beamerthememlsys}

\mlsyssetup{
  volume       = {Tutorial},
  chapter      = {Module 4},
  logo         = {../../../slides/assets/img/logo-mlsysbook.png},
  instlogo     = {../../../slides/assets/img/logo-harvard.png},
  chaptertitle = {MLSys·im: Design Space Exploration \& Synthesis},
}

% --- Fonts ---
\usepackage[T1]{fontenc}
\usepackage[scaled=0.9]{helvet}
\usepackage{courier}
\renewcommand{\familydefault}{\sfdefault}

% --- Packages ---
\usepackage{amsmath}
\usepackage{booktabs}
\usepackage[table]{xcolor}
\usepackage{listings}
\usepackage{tikz}
\usetikzlibrary{arrows.meta, positioning, calc, decorations.pathreplacing}

% --- Code listings ---
\lstset{
  language=Python,
  basicstyle=\ttfamily\scriptsize,
  keywordstyle=\color{crimson}\bfseries,
  stringstyle=\color{datastroke},
  commentstyle=\color{midgray}\itshape,
  backgroundcolor=\color{computeblue!20},
  frame=single,
  rulecolor=\color{computestroke},
  numbers=none,
  breaklines=true,
  columns=fullflexible,
  keepspaces=true,
  showstringspaces=false,
  xleftmargin=4pt,
  xrightmargin=4pt,
  aboveskip=6pt,
  belowskip=4pt,
}

% --- Convenience macros ---
\newcommand{\mlsysim}{\texttt{mlsysim}}
\newcommand{\wallbox}[2]{%
  \begin{block}{#1}#2\end{block}%
}
\newcommand{\PredictStart}{\begin{alertblock}{Predict Before You Peek}}
\newcommand{\PredictEnd}{\end{alertblock}}

% --- Image paths ---
\graphicspath{{images/}}

% --- Section count (must match actual \section{} count) ---
\setcounter{mlsystotalsections}{6}

\title{MLSys·im Tutorial --- Module 4}
\subtitle{Design Space Exploration \& Synthesis}
\author{Vijay Janapa Reddi}
\institute{Harvard University}
\date{Tutorial}

% =============================================================================
\begin{document}

% --- Roadmap ---
\begin{frame}[shrink]{Roadmap: Conference Tutorial}
\centering\small
\begin{tabular}{rll}
\toprule
\textbf{Module} & \textbf{Topic} & \textbf{Status} \\
\midrule
Module 1 & Foundations \& Architecture & \checkmark Done \\
Module 2 & Advanced Single-Node Analysis & \checkmark Done \\
Module 3 & Scale, Dollars, and Carbon & \checkmark Done \\
\rowcolor{crimson!12}
\textbf{Module 4} & \textbf{Design Space Exploration \& Synthesis} & \textbf{$\leftarrow$ You are here} \\
\bottomrule
\end{tabular}
\end{frame}


\section{Design Space Exploration}

% --- 6.1 Key Question ---
\begin{frame}[shrink]{Key Question}
\note{[1 min] Frame the search problem. The space is combinatorial.}
\begin{center}
\Large\bfseries
Given a budget and an SLA,\\[4pt]
how do you find the \emph{best} hardware configuration?
\end{center}
\vfill
\small
The search space is combinatorial:
$|\text{hardware}| \times |\text{batch sizes}| \times |\text{precisions}| \times |\text{parallelism configs}|$\\
can easily exceed $10^4$ configurations.
\end{frame}

% --- 6.2 The DSE Pattern ---
\begin{frame}[fragile,shrink=10]{The DSE Pattern: Declare, Search, Rank}
\note{[3 min] Three-step pattern. Emphasize: analytical models make exhaustive search feasible.}
\begin{columns}[T]
\column{0.5\textwidth}
\begin{enumerate}
  \item \textbf{Declare} the search space:
    \begin{itemize}\scriptsize
      \item Hardware: \{A100, H100, H200, B200\}
      \item Batch sizes: \{1, 2, 4, \ldots, 64\}
      \item Precisions: \{fp16, int8, int4\}
    \end{itemize}
  \item \textbf{Search} with an objective:
    \begin{itemize}\scriptsize
      \item \texttt{minimize: tco\_usd}
      \item \texttt{maximize: throughput}
    \end{itemize}
  \item \textbf{Rank} subject to constraints:
    \begin{itemize}\scriptsize
      \item \texttt{latency < 50 ms}
      \item \texttt{feasible == True}
    \end{itemize}
\end{enumerate}

\column{0.47\textwidth}
\begin{lstlisting}
from mlsysim.core.dse import DSE

dse = DSE(
  space={
    "batch_size": [1,4,8,16,32],
    "precision": ["fp16","int8"],
  },
  objective="maximize: throughput",
  constraints=["latency < 100"])
\end{lstlisting}
\vspace{0.3em}
\small
Analytical models $\Rightarrow$ each evaluation\\
takes $<$1\,ms $\Rightarrow$ exhaustive search.
\end{columns}
\end{frame}

% --- 6.3 Pareto Fronts ---
\begin{frame}[shrink]{Pareto Fronts: No Free Lunch}
\note{[2 min] Explain the Pareto front. The knee is usually the sweet spot.}
\centering
\includegraphics[width=0.65\textwidth]{images/pdf/pareto-front.pdf}

\vspace{0.5em}
\textbf{Pareto front}: the set of configurations where\\
improving one metric \emph{must} worsen another.

\begin{itemize}
  \item Lower latency $\Rightarrow$ lower throughput (smaller batch)
  \item Higher throughput $\Rightarrow$ higher cost (more GPUs)
  \item The ``knee'' of the curve is usually the sweet spot
\end{itemize}
\end{frame}

% --- 6.4 Live Demo: Engine.sweep ---
\begin{frame}[fragile,shrink=10]{Live Demo: Design Space Sweep}
\note{[3 min] Run live. Note multi-vendor hardware list: H100, MI300X, Gaudi3, B200.}
\begin{lstlisting}
hw_list = [mlsysim.Hardware.Cloud.H100,
           mlsysim.Hardware.Cloud.MI300X,
           mlsysim.Hardware.Cloud.Gaudi3,
           mlsysim.Hardware.Cloud.B200]
results = mlsysim.Engine.sweep(
    llama, hw_list, batch_sizes=[1, 4, 16, 64],
    precisions=["fp16", "int8"], efficiency=0.5)
for r in sorted(results,
        key=lambda x: x['profile'].latency.magnitude):
    p = r['profile']
    print(f"{r['hardware']:>8} bs={r['batch_size']:<3} "
          f"| {p.bottleneck:<8} | MFU={p.mfu:.3f}")
\end{lstlisting}
\end{frame}

% --- 6.5 Live Demo: DSE with Constraints ---
\begin{frame}[fragile,shrink=10]{Live Demo: DSE with Objective \& Constraints}
\note{[3 min] Run live. Show maximize throughput subject to latency constraint.}
\begin{lstlisting}
from mlsysim.core.dse import DSE

dse = DSE(
    space={"batch_size": [1,4,8,16,32,64],
           "precision": ["fp16", "int8"]},
    objective="maximize: throughput",
    constraints=["latency < 100"])
def evaluate(params):
    p = mlsysim.Engine.solve(llama, hw,
        batch_size=params["batch_size"],
        precision=params["precision"])
    return {"throughput": p.throughput.magnitude,
            "latency": p.latency.to("ms").magnitude}
best = dse.search(evaluate)
print(f"Best: {best['best_params']}")
\end{lstlisting}
\end{frame}

% --- 6.6 Batching Optimizer ---
\begin{frame}[fragile,shrink=10]{Live Demo: Batching Optimizer (Pareto Front)}
\note{[3 min] Run live. Show Pareto front of batch size vs latency.}
\begin{lstlisting}
from mlsysim.core.solver import BatchingOptimizer
opt = BatchingOptimizer()
result = opt.solve(
    model=llama, hardware=hw, seq_len=128,
    arrival_rate_qps=10.0,
    sla_latency_ms=20000.0, precision="fp16")
print(f"Optimal bs: {result.best_batch_size}")
for pt in result.pareto_front:
    print(f"  bs={pt['batch_size']:<4} "
          f"lat={pt['p99_latency'].m_as('ms'):.0f} ms")
\end{lstlisting}
\end{frame}

% --- 6.7 Exercise ---
\begin{frame}[fragile,shrink=10]{Exercise: Budget-Constrained Design}
\begin{alertblock}{Your Task (5 minutes)}
Given a \textbf{\$1M budget}, find the best hardware + batch size + precision
configuration for serving Llama-3-8B inference under a 50\,ms latency SLA.
\end{alertblock}

\begin{lstlisting}
# Hint: Use Engine.sweep across hardware tiers,
# filter by feasibility and latency SLA,
# then pick the highest-throughput config
# within budget (check unit_cost in Hardware)
\end{lstlisting}

\note{[5 min] Expected answer: H100 at INT8, batch size 16--32 gives the best
throughput/dollar under 50\,ms. A100 at FP16 fails the SLA at high batch sizes.
Turn to your neighbor: did you get the same answer? Why or why not? 60 seconds.}
\end{frame}

% --- 6.8 Takeaway ---
\begin{frame}[shrink]{Key Takeaway: Design Space Exploration}
\note{[1 min] Three points, repeat each.}
\begin{center}
\Large
\begin{enumerate}
  \item \textbf{Analytical models} make exhaustive DSE feasible ($<$1\,s for $10^4$ configs)
  \item \textbf{Pareto fronts} reveal the actual tradeoff structure
  \item \textbf{Constraints first}: filter infeasible, then optimize
\end{enumerate}

\vspace{1em}
\normalsize
\textit{``You cannot optimize what you cannot model.\\
You cannot model what you cannot measure.''}
\end{center}
\end{frame}


% =====================================================================
% PART 7: TINYML TO FRONTIER
% =====================================================================


\begin{frame}[fragile,shrink=10]{Live Demo: Programmatic Sweeps (Scenario A)}
\begin{lstlisting}[language=Python]
from mlsysim.core.engine import Engine
from mlsysim.hardware.registry import Hardware
from mlsysim.models.registry import Models
import pandas as pd

model = Models.Language.Llama3_8B
hardware = Hardware.Cloud.H100

results = []
for batch_size in [1, 8, 32, 128, 256]:
    perf = Engine.solve(model, hardware, batch_size=batch_size, precision="fp16")
    results.append({
        "Batch Size": batch_size,
        "Throughput (tok/s)": perf.throughput.magnitude,
        "Bottleneck": perf.bottleneck
    })

print(pd.DataFrame(results))
\end{lstlisting}
\end{frame}

\section{TinyML to Frontier}

% --- 7.1 Key Question ---
\begin{frame}[shrink]{Key Question}
\note{[1 min] ``Same equation, 9 orders of magnitude apart.''}
\begin{center}
\Large\bfseries
Can the same analytical framework model\\[4pt]
a \$2 microcontroller \emph{and} a \$3M GPU rack?
\end{center}
\vfill
\small
The answer is yes---because the \textbf{Roofline model} is universal.\\
Only the numbers change. The physics is the same.
\end{frame}

% --- 7.2 The Nine Orders of Magnitude ---
\begin{frame}[shrink]{The 9-Order-of-Magnitude Scale Span}
\note{[2 min] Walk through the table. Compute spans 10\^7x, power 10\^4.7x.}
\begin{columns}[T]
\column{0.52\textwidth}
\centering
\includegraphics[width=0.85\textwidth]{images/pdf/hardware-spectrum.pdf}

\column{0.45\textwidth}
\scriptsize
\begin{tabular}{lrr}
\toprule
\textbf{Device} & \textbf{FLOPS} & \textbf{TDP} \\
\midrule
nRF52840  & 64\,M  & 15\,mW \\
ESP32-S3  & 500\,M & 400\,mW \\
\rowcolor{gray!15}
H100 SXM  & 989\,T & 700\,W \\
\bottomrule
\end{tabular}

\vspace{0.5em}
\begin{tabular}{lr}
\textbf{Compute} & $\sim 10^{7}\times$ \\
\textbf{Mem BW}  & $\sim 10^{7}\times$ \\
\textbf{Power}   & $\sim 10^{4.7}\times$ \\
\end{tabular}
\end{columns}

\vspace{0.5em}
\centering
\alert{Same Roofline equation. Nine orders of magnitude apart.}
\end{frame}

% --- 7.3 TinyML Memory Hierarchy ---
\begin{frame}[shrink]{Flash vs SRAM: The TinyML Memory Wall}
\note{[3 min] Key difference: TinyML has Flash (8 MB, 80 MB/s) vs Cloud HBM (80 GB, 3.35 TB/s).}
\begin{columns}[T]
\column{0.5\textwidth}
\textbf{Cloud GPU (H100)}
\begin{itemize}
  \item Weights in HBM (80\,GB)
  \item Activations in SRAM (50\,MB L2)
  \item Bandwidth: 3.35\,TB/s
  \item Model $\leq$ 80\,GB $\Rightarrow$ fits
\end{itemize}

\column{0.5\textwidth}
\textbf{TinyML MCU (ESP32-S3)}
\begin{itemize}
  \item Weights in \textbf{Flash} (8\,MB)
  \item Activations in \textbf{SRAM} (512\,KiB)
  \item Flash BW: 80\,MB/s
  \item Model $\leq$ 8\,MB $\Rightarrow$ fits
  \item Model $\leq$ 512\,KiB $\Rightarrow$ SRAM-only (12$\times$ faster)
\end{itemize}
\end{columns}

\vspace{1em}
\textbf{Key insight:} mlsysim automatically selects the right memory tier:
\begin{enumerate}
  \item Model fits in SRAM $\Rightarrow$ use SRAM bandwidth
  \item Model in Flash $\Rightarrow$ use Flash bandwidth (bottleneck!)
  \item Model exceeds Flash $\Rightarrow$ infeasible
\end{enumerate}
\end{frame}

% --- 7.4 Energy per Inference ---
\begin{frame}[shrink]{Energy per Inference: $\mu$J to Joules}
\note{[2 min] 6 orders of magnitude in energy. At TinyML scale, battery life is the constraint.}
\centering
\begin{tabular}{lrrl}
\toprule
\textbf{Device} & \textbf{TDP} & \textbf{Energy/Inf} & \textbf{Use Case} \\
\midrule
nRF52840       & 15\,mW  & $\sim$50\,$\mu$J & Keyword spotting \\
ESP32-S3       & 400\,mW & $\sim$1\,mJ      & Person detection \\
Jetson Orin NX & 25\,W   & $\sim$25\,mJ     & Object detection \\
\rowcolor{gray!15}
H100 SXM       & 700\,W  & $\sim$50\,J      & LLM inference \\
\bottomrule
\end{tabular}

\vspace{1em}
\begin{itemize}
  \item 6 orders of magnitude in energy per inference
  \item TinyML: battery life is the primary constraint
  \item Duty cycling (sleep 99\% of the time) extends battery to years
\end{itemize}
\end{frame}

% --- 7.5 Live Demo: Hardware Comparison ---
\begin{frame}[fragile,shrink=10]{Live Demo: nRF52840 vs ESP32 vs H100}
\note{[3 min] Run live. nRF52840 is memory-bound (Flash). H100 finishes in microseconds.}
\begin{lstlisting}
tiny_model = mlsysim.Models.Tiny.KeywordSpotting
devices = [mlsysim.Hardware.Tiny.nRF52840,
           mlsysim.Hardware.Tiny.ESP32_S3,
           mlsysim.Hardware.Cloud.H100]
for d in devices:
    p = mlsysim.Engine.solve(tiny_model, d,
          precision="int8", efficiency=0.3)
    print(f"{d.name:>15}: {p.latency.to('ms'):.2f~P}"
          f" | {p.bottleneck:<8}")
\end{lstlisting}

\vspace{0.3em}
\begin{exampleblock}{Observation}
nRF52840 is memory-bound (Flash bottleneck).
H100 finishes in $\mu$s---but wastes 700\,W doing it.
\textbf{Right-sizing hardware matters.}
mlsysim also supports \textbf{Coral Edge TPU}, Jetson Orin, and Inferentia2.
\end{exampleblock}
\end{frame}

% --- 7.6 Same Roofline, Different Physics ---
\begin{frame}[fragile,shrink=10]{Same Roofline, Different Physics}
\note{[2 min] Both are memory-bound at batch size 1. Same equation, same diagnosis.}
\begin{columns}[T]
\column{0.5\textwidth}
\begin{center}
\[
\text{Latency} = \max\!\left(
  \frac{\text{FLOPs}}{\text{Peak} \times \eta},\;
  \frac{\text{Weights}}{\text{BW}}
\right)
\]
\end{center}

\vspace{0.3em}
\scriptsize
\begin{tabular}{lll}
\toprule
\textbf{Term} & \textbf{Cloud} & \textbf{TinyML} \\
\midrule
Peak      & 989\,TF & 500\,MF \\
BW        & 3.35\,TB/s & 80\,MB/s \\
Overhead  & 0.01\,ms & 1.0\,ms \\
\bottomrule
\end{tabular}

\column{0.47\textwidth}
\begin{lstlisting}
# Same API, different hardware
for d in [mlsysim.Hardware.Cloud.H100,
          mlsysim.Hardware.Tiny.ESP32_S3]:
    p = mlsysim.Engine.solve(
        mlsysim.Models.Tiny
          .KeywordSpotting,
        d, precision="int8")
    print(f"{d.name}: {p.bottleneck}")
\end{lstlisting}
\end{columns}

\vspace{0.3em}
\centering
\small
\textbf{Both are memory-bound at batch size 1.} Same equation, same diagnosis.
\end{frame}

% --- 7.7 Exercise ---
\begin{frame}[fragile,shrink=10]{Exercise: Right-Sizing for IoT}
\begin{alertblock}{Your Task (5 minutes)}
A smart doorbell needs person detection at $<$100\,ms latency on a
coin-cell battery (500\,mAh @ 3V = 1.5\,Wh). Compare ESP32-S3 vs
nRF52840: which device can run inference for a full year at 1 inference/minute?
\end{alertblock}

\begin{lstlisting}
import mlsysim

model = mlsysim.Models.Tiny.PersonDetection
for hw in [mlsysim.Hardware.Tiny.ESP32_S3,
           mlsysim.Hardware.Tiny.nRF52840]:
    p = mlsysim.Engine.solve(model, hw,
          precision="int8", efficiency=0.3)
    energy_j = p.energy.to("J").magnitude
    inferences_per_year = 60 * 24 * 365
    # YOUR CODE: total energy vs 1.5 Wh battery
\end{lstlisting}

\note{[5 min] Expected: nRF52840 at 15\,mW lasts $>$1 year with duty cycling.
ESP32-S3 at 400\,mW drains the battery in weeks.
Turn to your neighbor: did you get the same answer? Why or why not? 60 seconds.}
\end{frame}

% --- 7.8 Takeaway ---
\begin{frame}[shrink]{Key Takeaway: TinyML to Frontier}
\note{[1 min] Three points. Right-size hardware to the workload.}
\begin{center}
\Large
\begin{enumerate}
  \item The Roofline model spans \textbf{9 orders of magnitude}
  \item \textbf{Right-size} hardware to the workload, not the other way around
  \item At TinyML scale, \textbf{energy and memory} dominate; at cloud scale, \textbf{cost and communication}
\end{enumerate}

\vspace{1em}
\normalsize
\textit{``The best accelerator is the smallest one\\
that meets your latency and accuracy SLA.''}
\end{center}
\end{frame}


% =====================================================================
% PART 8: ADVANCED TOPICS
% =====================================================================

\section{Advanced Topics}

% --- 8.1 Key Question ---
\begin{frame}[shrink]{Key Question}
\note{[1 min] Single-wall analysis is not enough. You need the full pipeline.}
\begin{center}
\Large\bfseries
How do you compose multiple analytical models\\[4pt]
into an end-to-end system analysis?
\end{center}
\vfill
\small
Single-wall analysis answers ``what is the bottleneck?''\\
\textbf{Pipeline composition} answers ``what is the total cost of the whole stack?''
\end{frame}

% --- 8.2 Pipeline Composition Pattern ---
\begin{frame}[fragile,shrink=10]{The Pipeline Composition Pattern}
\note{[3 min] Show the Pipeline pattern. Each stage feeds the next.}
\begin{lstlisting}
from mlsysim.core.pipeline import Pipeline
from mlsysim.core.solver import (
    ScalingModel, DistributedModel, EconomicsModel)

# Chain: Algorithm -> Fleet -> Economics
pipe = Pipeline([
    ScalingModel(),      # Wall 11: compute budget
    DistributedModel(),  # Wall 14: comms overhead
    EconomicsModel(),    # Wall 17: total cost
])
print(pipe.explain())  # shows DAG + gaps
\end{lstlisting}

\vspace{0.5em}
\texttt{explain()} shows:
\begin{itemize}
  \item What each stage requires and produces
  \item Which walls from the 22-wall taxonomy are covered
  \item Where gaps exist (missing inputs)
\end{itemize}
\end{frame}

% --- 8.3 Pipeline Run ---
\begin{frame}[fragile,shrink=10]{Running the Pipeline}
\note{[3 min] Run live. Show all stages and their outputs.}
\begin{lstlisting}
from mlsysim.core.constants import Q_

results = pipe.run(
    compute_budget=Q_("1e21 FLOP"),
    fleet=mlsysim.Systems.Clusters.Research_256,
    duration_days=30,
    datacenter=mlsysim.Infrastructure.Grids.Quebec,
    mfu=0.45,
    infrastructure_multiplier=2.5)

for stage_name, result in results.items():
    print(f"\n--- {stage_name} ---")
    print(result)
\end{lstlisting}

\vspace{0.5em}
\textbf{Key design principle:} Each resolver's output fields become available
as inputs to subsequent stages. The pipeline is \emph{not} a black box---it
is a transparent analytical tool.
\end{frame}

% --- 8.4 Inference-Time Compute (The Reasoning Wall) ---
\begin{frame}[fragile,shrink=10]{Wall 12: Inference-Time Compute (The Reasoning Wall)}
\note{[3 min] "With models like OpenAI o1, compute scaling shifts from training to inference. The model generates K hidden reasoning steps before answering."}
\begin{lstlisting}
from mlsysim.core.solver import InferenceScalingModel

reasoning_solver = InferenceScalingModel()
result = reasoning_solver.solve(
    model=llama, hardware=hw,
    reasoning_steps=128,  # K hidden CoT steps
    precision="fp16", efficiency=0.5)

print(f"Reasoning Time: {result.total_reasoning_time.to('s'):.1f}")
print(f"Energy per Query: {result.energy_per_query.to('J'):.1f}")
\end{lstlisting}

\vspace{0.3em}
\small
\textbf{The Paradigm Shift:}
\begin{itemize}
  \item Standard LLM: Latency scales with output token count.
  \item Reasoning LLM (o1-like): Latency scales with $K$ hidden Chain-of-Thought steps.
  \item \textbf{Impact:} Inference shifts back toward being \emph{compute-bound} as generation length dominates prefill.
\end{itemize}
\end{frame}

% --- 8.5 Sensitivity Analysis ---
\begin{frame}[fragile,shrink=10]{Wall 21: Sensitivity Analysis}
\note{[3 min] ``Which knob should I turn next?'' The parameter with the largest partial derivative.}
\begin{lstlisting}
from mlsysim.core.solver import SensitivitySolver
solver = SensitivitySolver()
result = solver.solve(
    model=llama, hardware=hw, precision="fp16",
    perturbation_pct=10.0, efficiency=0.5)
print(f"Binding: {result.binding_constraint}")
for p, s in result.sensitivities.items():
    tag = "<<<" if p == result.binding_constraint else ""
    print(f"  {p:>20}: {s:+.4f}  {tag}")
\end{lstlisting}

\vspace{0.3em}
\small
\textbf{Rule:} Invest in the parameter with the \emph{largest} $|\partial T / \partial p|$.\\
Improving a non-binding parameter yields \textbf{zero} measurable gain.
\end{frame}

% --- 8.6 Inverse Roofline (Synthesis) ---
\begin{frame}[fragile,shrink=10]{Wall 22: Inverse Roofline (SynthesisSolver)}
\note{[3 min] Most powerful move: derive hardware requirements from SLA. ``I need 50ms. What hardware?''}
\begin{lstlisting}
from mlsysim.core.constants import Q_
from mlsysim.core.solver import SynthesisSolver

solver = SynthesisSolver()
result = solver.solve(
    model=llama, target_latency=Q_("50 ms"),
    precision="fp16", efficiency=0.5)
print(f"Required BW:    {result.required_bw:.0f~P}")
print(f"Required FLOPS: {result.required_flops:.1f~P}")
\end{lstlisting}

\vspace{0.5em}
\textbf{Usage:} ``I need 50\,ms TTFT. What hardware do I need?''\\
$\Rightarrow$ Derive specs from SLA, then match to real accelerators.
\end{frame}

% --- 8.7 Fallacies & Pitfalls ---
\begin{frame}[shrink]{Fallacies \& Pitfalls (Patterson Tradition)}
\note{[2 min] Walk through each fallacy with quantitative rebuttal.}
\begin{block}{Fallacy: ``Cheaper hardware is always more cost-effective''}
Reality: A 2$\times$ cheaper GPU that takes 3$\times$ longer has \emph{higher} TCO.\\
$\text{TCO} = \text{CapEx} + \text{OpEx}$; slower hardware burns more electricity.
\end{block}

\begin{block}{Fallacy: ``More FLOPS = proportionally faster''}
Reality: A100 $\to$ H100 is 6.3$\times$ peak FLOPS but only $\sim$2$\times$ for
memory-bound LLM inference. The binding constraint is bandwidth, not compute.
\end{block}

\begin{block}{Pitfall: Optimizing a non-binding parameter}
If inference is memory-bound, doubling FLOPS gives 0\% speedup.\\
Use \texttt{SensitivitySolver} to find the binding constraint \emph{first}.
\end{block}

\begin{block}{Pitfall: Ignoring the infrastructure multiplier}
GPU cost is 40\% of total datacenter TCO. Budgeting for GPUs alone\\
underestimates the true cost by 2--3$\times$.
\end{block}
\end{frame}

% --- 8.8 The Iron Law Revisited ---
\begin{frame}[shrink]{The Iron Law Revisited: All Five Terms}
\note{[2 min] Return to the master equation. Every wall maps to one of these five terms.}
\begin{center}
\Large
\[
T = \frac{\text{FLOPs}}{N \times \text{Peak} \times \text{MFU} \times \eta_{\text{scale}} \times \text{Goodput}}
\]
\end{center}

\vspace{0.5em}
\centering
\begin{tabular}{lll}
\toprule
\textbf{Term} & \textbf{What reduces it} & \textbf{Wall} \\
\midrule
$N$                  & Budget (buy more GPUs)        & Wall 17 \\
Peak                 & GPU generation (H100$\to$B200)& Wall 1 \\
MFU                  & FlashAttention, kernel fusion  & Wall 3 \\
$\eta_{\text{scale}}$& Network BW, gradient compression & Wall 14 \\
Goodput              & Checkpointing, fault tolerance & Walls 15, 19 \\
\bottomrule
\end{tabular}

\vspace{0.5em}
\textbf{Every wall in the 22-wall taxonomy maps to one of these five terms.}\\
That is the entire framework.
\end{frame}

% --- 8.9 Exercise ---
\begin{frame}[fragile,shrink=10]{Exercise: Binding Constraint Analysis}
\note{[5 min] At bs=1, memory bandwidth is binding. At bs=64, peak FLOPS. This IS the Roofline transition.
Turn to your neighbor: did you get the same answer? Why or why not? 60 seconds.}
\begin{alertblock}{Your Task (5 minutes)}
Use \texttt{SensitivitySolver} on Llama-3-8B + H100.
Which parameter is binding? Now switch to \texttt{batch\_size=64} and
re-run. Does the binding constraint change?
\end{alertblock}

\begin{lstlisting}
from mlsysim.core.solver import SensitivitySolver
solver = SensitivitySolver()
# Batch size 1
r1 = solver.solve(llama, hw, efficiency=0.5)
print(f"bs=1:  binding = {r1.binding_constraint}")
# Batch size 64
p64 = mlsysim.Engine.solve(llama, hw, batch_size=64)
print(f"bs=64: bottleneck = {p64.bottleneck}")
\end{lstlisting}

\end{frame}


% =====================================================================
% PART 9: WRAP-UP & FUTURE
% =====================================================================


\begin{frame}[fragile,shrink=10]{Live Demo: Inverting the Roofline (Scenario D)}
Instead of asking "How fast is this GPU?", what if we ask "What hardware do I need to buy to meet my 30ms latency SLA?"
\begin{lstlisting}[language=Python]
from mlsysim.solvers import SynthesisSolver
from mlsysim.models.registry import Models
from mlsysim.core.constants import Q_

solver = SynthesisSolver()
requirements = solver.solve(
    model=Models.Language.Llama3_8B,
    target_latency=Q_("30 ms"),
    batch_size=1,
    precision="fp16"
)

print(f"Required HBM Bandwidth: {requirements.required_bw.to('GB/s'):.1f}")
print(f"Required Compute:       {requirements.required_flops.to('TFLOPs/s'):.1f}")
\end{lstlisting}
\end{frame}


\section{Extensibility \& Transparency}
\begin{frame}[shrink]{Not a Black Box: The C.A.T. Framework}
\mlsysim{} is designed as an executable textbook using the \textbf{Concept, Algebra, Tool (C.A.T)} framework.
\vfill
\begin{itemize}
  \item \textbf{Transparent Math:} Every solver is a pure Python function implementing published academic algebra. You can read the code to see exactly how $T_{\text{comm}}$ or ITL is calculated.
  \item \textbf{Extensible:} If you want to model a novel MoE routing algorithm or a new network topology, you don't need to rewrite a C++ simulator.
  \item \textbf{How:} Simply subclass \texttt{BaseModel} or \texttt{BaseOptimizer} and write the Python math. It automatically integrates into the \texttt{Pipeline} architecture.
\end{itemize}
\end{frame}

\section{Wrap-up \& Future}

% --- 9.1 The 22 Walls in One Slide ---
\begin{frame}[shrink]{The 22 Walls of ML Systems}
\note{[2 min] Full recap. Every wall maps to one Iron Law term.}
\centering
\small
\begin{tabular}{clcl}
\toprule
\textbf{\#} & \textbf{Wall} & \textbf{\#} & \textbf{Wall} \\
\midrule
1  & Compute         & 12 & Reasoning (Emerging) \\
2  & Memory          & 13 & Fidelity (Compression) \\
3  & Software (MFU)  & 14 & Communication \\
4  & Serving         & 15 & Fragility (Reliability) \\
5  & Batching (KV)   & 16 & Multi-tenant (Queueing) \\
6  & Streaming       & 17 & Capital (TCO) \\
7  & Tail Latency    & 18 & Sustainability \\
8  & Ingestion       & 19 & Checkpoint \\
9  & Transformation  & 20 & Safety (Privacy) \\
10 & Locality        & 21 & Sensitivity \\
11 & Complexity      & 22 & Synthesis (Inverse) \\
\bottomrule
\end{tabular}

\vspace{0.5em}
\textbf{6 Domains:} Node $\cdot$ Data $\cdot$ Algorithm $\cdot$ Fleet $\cdot$ Operations $\cdot$ Meta-Analysis

\vspace{0.3em}
\alert{Every wall maps to one term in the Iron Law.}
\end{frame}

% --- 9.2 What This Tool Does NOT Model ---
\begin{frame}[shrink]{What \texttt{mlsysim} Does \emph{Not} Model}
\note{[2 min] Honest about limitations. Not a simulator, not a profiler. A reasoning framework.}
\begin{columns}[T]
\column{0.5\textwidth}
\textbf{Not modeled (v0.1.0):}
\begin{itemize}
  \item Cache effects / tiling / fusion
  \item Real network congestion / jitter
  \item CUDA kernel scheduling
  \item Multi-model co-location
  \item Compiler optimizations
\end{itemize}

\column{0.5\textwidth}
\textbf{By design:}
\begin{itemize}
  \item Not a cycle-accurate simulator
  \item Not a profiler replacement
  \item Not a deployment tool
  \item \textbf{Is:} a first-principles\\
        analytical reasoning tool
\end{itemize}

\vspace{0.5em}
\textbf{Use it for:}
\begin{itemize}
  \item ``Which constraint is binding?''
  \item ``Is this hardware sufficient?''
  \item ``What should I try first?''
\end{itemize}
\end{columns}

\vspace{0.5em}
\centering
\small\itshape
``All models are wrong, but some are useful.'' --- George Box
\end{frame}

% --- 9.3 v0.2.0 Roadmap ---
\begin{frame}[shrink]{v0.2.0 Roadmap}
\note{[2 min] Quick overview. Contributions welcome.}
\small
\begin{columns}[T]
\begin{column}{0.48\textwidth}
\textbf{Models \& Hardware}
\begin{enumerate}
  \item \textbf{MoE support} --- Expert routing + parallelism
  \item \textbf{GQA / MQA} --- Grouped-query KV cache
  \item \textbf{Network congestion} --- Contention-aware collectives
\end{enumerate}
\end{column}
\begin{column}{0.48\textwidth}
\textbf{Tooling \& Integration}
\begin{enumerate}\setcounter{enumi}{3}
  \item \textbf{Spot pricing} --- Cloud cost optimization
  \item \textbf{Browser exercises} --- Live exploration
  \item \textbf{HuggingFace Hub} --- Auto-import models
\end{enumerate}
\end{column}
\end{columns}

\vspace{0.5em}
\begin{center}
Contributions welcome: \texttt{github.com/harvard-edge/mlsysim}
\end{center}
\end{frame}

% --- 9.4 Design Challenge (Capstone) ---
\begin{frame}[fragile,shrink=10]{Design Challenge: Capstone Exercise}
\begin{alertblock}{The Problem}
\textbf{\$5M budget.} Serve Llama-3 70B at \textbf{1{,}000 QPS} with
\textbf{$<$100\,ms TTFT} in \textbf{two regions} (US-East + EU-West).
Design the fleet.
\end{alertblock}

\vspace{0.3em}
\begin{columns}[T]
\column{0.48\textwidth}
\textbf{You must specify:}
\begin{enumerate}\scriptsize
  \item Hardware choice (which GPU? how many?)
  \item Parallelism strategy (TP $\times$ PP)
  \item Precision (FP16? INT8? FP8?)
  \item Geographic placement (carbon?)
  \item Redundancy (replicas per region?)
\end{enumerate}

\column{0.49\textwidth}
\begin{lstlisting}
import mlsysim

# Capstone starter code
model = mlsysim.Models.Language.Llama3_70B
hw = mlsysim.Hardware.Cloud.H100

# Step 1: Can it fit?
p = mlsysim.Engine.solve(
    model, hw, precision="fp16")
print(f"Feasible: {p.feasible}")
print(f"Latency:  {p.latency:~P}")

# Step 2: TCO for your fleet
econ = mlsysim.EconomicsModel()
# YOUR CODE: design the fleet
\end{lstlisting}
\end{columns}

\note{[35 min] This is intentionally under-specified. Students must make and defend
assumptions. Key insight: 70B at FP16 = 140\,GB, needs tensor parallelism
across 2+ H100s per replica. At 1000 QPS with 100\,ms TTFT, you need many
replicas. Budget constrains how many.}
\end{frame}

% --- 9.4b: Personal Transfer Moment ---
\begin{frame}[shrink]{Your Turn: Name That Wall}
\note{[3 min] THE TRANSFER MOMENT. Silent writing. Then 2-3 volunteers share.}
\centering
\Large
\textbf{Think of one ML system in your own work.}\\[1cm]
\normalsize
Write down:\\[0.3cm]
\begin{enumerate}
  \item What is the system? (model + hardware + use case)
  \item Which of the 22 walls is the \textbf{binding constraint}?
  \item What would you change to move the wall?
\end{enumerate}

\vfill
\small\textcolor{gray}{2 minutes. Then we will hear from 3 volunteers.}
\end{frame}

% --- 9.5 Resources & Q&A ---
\begin{frame}[shrink]{Resources \& Next Steps}
\note{[3 min] Point to: mlsysim GitHub, textbook, cheatsheet. ``Take a photo of the cheatsheet.''}
\begin{columns}[T]
\column{0.55\textwidth}
\textbf{Get Started}
\begin{itemize}
  \item \texttt{pip install mlsysim}
  \item GitHub: \texttt{harvard-edge/mlsysim}
  \item \textbf{Code Cookbook:} See the 5 interactive scenarios in the Appendix!
  \item Full docs: \texttt{mlsysim.readthedocs.io}
\end{itemize}

\vspace{0.5em}
\textbf{The Textbook}
\begin{itemize}
  \item \emph{Machine Learning Systems}
  \item Volume I: Foundations (single node)
  \item Volume II: Systems at Scale (fleet)
  \item \texttt{mlsysbook.ai}
\end{itemize}

\column{0.40\textwidth}
\textbf{Key Papers}
\begin{itemize}
  \item Williams et al.\ (2009)\\
        {\scriptsize Roofline Model}
  \item Chowdhery et al.\ (2022)\\
        {\scriptsize PaLM / MFU}
  \item Hoffmann et al.\ (2022)\\
        {\scriptsize Chinchilla Scaling}
  \item Patterson et al.\ (2021)\\
        {\scriptsize Carbon \& Training}
  \item OpenAI (2024)\\
        {\scriptsize o1 / Reasoning Scaling}
\end{itemize}
\end{columns}

\vspace{1em}
\begin{center}
\Large\bfseries Thank you! Questions?
\end{center}
\end{frame}



\end{document}
